\documentclass[11pt,a4paper]{article}
\usepackage[utf8]{inputenc}
\usepackage[T1]{fontenc}
\usepackage[french]{babel}
\usepackage{geometry}
\usepackage{tikz}
\usetikzlibrary{arrows.meta,positioning}
\geometry{margin=2.2cm}

\title{Rapport technique court\\Architecture et infrastructure}
\author{}
\date{}

\begin{document}
\maketitle

\section{Introduction}
Ce projet est un backend pour une plateforme e-commerce proche d'un r\'eseau social. 
Il g\`ere des fonctions logistiques importantes: paiements, livraisons, catalogue produit et stock.
L'architecture suit un mod\`ele microservices: chaque service est s\'epar\'e selon une responsabilit\'e m\'etier claire.

\section{Architecture et infrastructure}

\subsection{1. APIs REST}
Chaque microservice expose une API REST.
Cette API permet aux autres services et aux clients d'appeler les fonctionnalit\'es de mani\`ere standard (HTTP + JSON).
Le d\'ecoupage par service garde le code plus simple \`a maintenir et \`a faire \'evoluer.

\subsection{2. Authentification}
L'authentification repose sur des tokens JWT.
Un token contient des \textit{claims} (par exemple l'identit\'e utilisateur, le r\^ole, la dur\'ee de validit\'e).
Le syst\`eme utilise:
\begin{itemize}
  \item un \textbf{access token} (courte dur\'ee) pour acc\'eder aux APIs;
  \item un \textbf{refresh token} (dur\'ee plus longue) pour obtenir un nouvel access token.
\end{itemize}

La signature est v\'erifi\'ee localement dans chaque service avec des cl\'es RS256.
Donc, un service n'a pas besoin d'appeler le service d'auth \`a chaque requ\^ete pour valider un token.

\subsection{3. Coh\'erence d'\'etat avec RabbitMQ}
Les services communiquent de mani\`ere asynchrone via RabbitMQ pour garder une coh\'erence globale.
Le r\^ole utilisateur est un bon exemple:
\begin{itemize}
  \item un utilisateur cr\'ee une boutique dans \texttt{store\_service};
  \item un \'ev\'enement est publi\'e dans RabbitMQ;
  \item \texttt{user\_service} consomme l'\'ev\'enement et passe le r\^ole de \texttt{BUYER} \`a \texttt{SELLER}.
\end{itemize}

Exception importante: la validation de paiement demande un traitement synchronis\'e (r\'eponse imm\'ediate attendue).
En pratique, chaque service ex\'ecute au minimum:
\begin{itemize}
  \item un serveur REST;
  \item un processus consommateur (\texttt{start\_consumer}) pour les messages RabbitMQ.
\end{itemize}

\subsection{4. Discovery et registry}
Le projet utilise Consul pour la d\'ecouverte de services.
Par exemple, dans \texttt{services/delivery\_service/delivery/consul\_registry.py}, le service s'enregistre avec ses m\'etadonn\'ees r\'eseau.
Les valeurs d'adresse et de port proviennent du lancement (voir \texttt{services/delivery\_service/run.sh}).
Ainsi, les autres services peuvent trouver dynamiquement \texttt{delivery\_service} sans configuration statique.

\section{Sch\'ema global}
\begin{center}
\begin{tikzpicture}[
  node distance=1.6cm and 2cm,
  box/.style={draw, rounded corners, align=center, minimum width=2.8cm, minimum height=0.9cm},
  broker/.style={draw, rounded corners, fill=gray!10, align=center, minimum width=3.2cm, minimum height=0.9cm},
  >={Stealth}
]
\node[box] (client) {Client/API Gateway};
\node[box, below left=of client] (store) {store\_service\\REST + consumer};
\node[box, below=of client] (user) {user\_service\\REST + consumer};
\node[box, below right=of client] (delivery) {delivery\_service\\REST + consumer};
\node[box, right=3.3cm of client] (auth) {auth\_service\\JWT RS256};
\node[broker, below=2.4cm of user] (mq) {RabbitMQ\\\'ev\'enements m\'etier};
\node[box, right=3.3cm of mq] (consul) {Consul\\service registry};

\draw[->] (client) -- (store);
\draw[->] (client) -- (user);
\draw[->] (client) -- (delivery);
\draw[->] (client) -- (auth);

\draw[->] (store) -- (mq);
\draw[->] (mq) -- (user);
\draw[->] (delivery) -- (mq);

\draw[->, dashed] (store) -- (consul);
\draw[->, dashed] (user) -- (consul);
\draw[->, dashed] (delivery) -- (consul);
\end{tikzpicture}
\end{center}

\end{document}
